\documentclass[12pt,a4paper]{article}

\usepackage[utf8]{inputenc}
\usepackage[T1]{fontenc}
\usepackage{geometry}
\usepackage{amsmath,amssymb}
\usepackage{longtable}
\usepackage{array}

\geometry{margin=2.5cm}
\newcommand{\instituicao}{Universidade Federal Fluminense}
\newcommand{\unidade}{Instituto de Física}
\newcommand{\disciplina}{Física Experimental IV}

\begin{document}

\begin{titlepage}
  \centering
  {\Large \unidade \par}
  \vspace{0.4cm}
  {\large \instituicao \par}
  \vfill
  {\LARGE\bfseries \disciplina \par}
  \vspace{0.8cm}
  {\Huge\bfseries Prática 1 -- Revisão \par}
  \vspace{0.8cm}
  {\large Roteiro em \LaTeX\ estruturado a partir do conteúdo do portal da disciplina.\par}
  \vfill
  {\large 13 de julho de 2026 \par}
\end{titlepage}

\tableofcontents
\newpage

\section{Objetivos de aprendizagem}

\begin{itemize}
  \item Rever o uso correto do multímetro digital em medidas de resistência, tensão e corrente.
  \item Escrever a incerteza de uma medida direta usando a especificação do aparelho.
  \item Calcular a resistência equivalente de resistores em série e propagar as incertezas por dois procedimentos.
  \item Comparar uma corrente calculada com uma corrente medida diretamente no circuito.
  \item Reconhecer que a presença do LED exige cuidado na interpretação da queda de tensão total do circuito.
\end{itemize}

\section{Introdução teórica}

Esta prática revisa dois pontos: uso do multímetro em medidas diretas e escrita correta de incertezas experimentais. Antes de medir,
defina a grandeza de interesse, escolha a função correta, a escala adequada e o borne apropriado. Resistências e tensões são medidas
em paralelo ao elemento de interesse. Correntes são medidas em série, com o circuito aberto no ponto de medição. Como o circuito
contém três resistores em série com um LED, a tensão da fonte não deve ser identificada automaticamente com a soma das quedas de
tensão apenas nos resistores.

\section{Medidas diretas e incerteza do multímetro}

A leitura do multímetro fornece o valor medido. A incerteza depende da faixa utilizada, da resolução e da especificação do aparelho
no manual. Uma forma usual de estimá-la é:

\begin{equation}
\Delta = (\text{erro percentual}) \cdot (\text{valor lido}) + (\text{numero de digitos}) \cdot (\text{resolucao})
\end{equation}

O primeiro termo depende do valor lido. O segundo depende da resolução da faixa utilizada.

Para usar essa expressão, identifique no manual:

\begin{itemize}
  \item a especificação percentual da faixa escolhida;
  \item o número de dígitos que deve ser somado ao cálculo;
  \item a resolução correspondente à menor variação exibida na tela.
\end{itemize}

Ao final, a medida deve ser escrita com:

\begin{itemize}
  \item no máximo dois algarismos significativos na incerteza;
  \item o valor central e a incerteza com o mesmo número de casas decimais;
  \item a unidade fora do parêntese.
\end{itemize}

\section{Propagação de incertezas}

Quando uma grandeza é calculada a partir de outras medidas, sua incerteza deve ser obtida a partir das incertezas das variáveis
medidas.

A forma linear soma os módulos das contribuições individuais e fornece uma estimativa conservadora:

\begin{equation}
\Delta f =
\left| \frac{\partial f}{\partial x} \right| \Delta x +
\left| \frac{\partial f}{\partial y} \right| \Delta y +
\left| \frac{\partial f}{\partial z} \right| \Delta z + \cdots
\end{equation}

A forma em quadratura combina as contribuições de modo estatístico:

\begin{equation}
\Delta f =
\sqrt{
\left( \frac{\partial f}{\partial x} \right)^2 (\Delta x)^2 +
\left( \frac{\partial f}{\partial y} \right)^2 (\Delta y)^2 +
\left( \frac{\partial f}{\partial z} \right)^2 (\Delta z)^2 + \cdots
}
\end{equation}

Nesta prática, as duas expressões serão aplicadas à resistência equivalente e à corrente calculada.

\section{Descrição do experimento}

O experimento usa uma fonte DC ajustável e um circuito com três resistores em série com um LED. Medem-se as resistências, as quedas
de tensão nos resistores e a corrente do circuito. Depois, calcula-se a resistência equivalente e compara-se a corrente calculada com
a corrente medida.

\section{Materiais e montagem}

\begin{itemize}
  \item 01 fonte de tensão DC ajustável;
  \item 01 kit de circuito com três resistores e 1 LED;
  \item 01 multímetro digital com manual;
  \item 02 pontas de prova;
  \item 02 fios com terminais banana.
\end{itemize}

A montagem principal é um circuito em série contendo três resistores e um LED, alimentado por uma fonte DC ajustada para 10 V.

\section{Procedimento experimental}

\begin{enumerate}
  \item Antes de qualquer conexão, identifique se a próxima medida será de resistência, tensão ou corrente.
  \item Selecione a função apropriada no multímetro e, sempre que necessário, comece na maior escala.
  \item Conecte uma ponta de prova ao borne comum e a outra ao borne adequado para a grandeza escolhida.
  \item Meça separadamente as três resistências do kit e registre cada valor com sua incerteza.
  \item Monte o circuito com os três resistores em série com o LED e ajuste a fonte para 10 V.
  \item Meça a tensão em cada resistor, sempre em paralelo ao elemento observado, e registre o resultado com a incerteza.
  \item Calcule a resistência equivalente dos três resistores e propague a incerteza pelos dois métodos do roteiro.
  \item Estime a corrente do circuito e propague sua incerteza pelos dois métodos.
  \item Meça diretamente a corrente no circuito, agora com o multímetro em série, e compare com o valor calculado.
  \item Discuta a compatibilidade entre os resultados e o papel do LED na interpretação do circuito.
\end{enumerate}

\paragraph{Cuidado operacional.}
Medidas de resistência devem ser feitas sem o circuito alimentado. Medidas de corrente exigem abrir o circuito e inserir o multímetro
em série. Além disso, não aplique a tensão da fonte diretamente aos terminais do LED sem resistor em série: a tensão direta
admissível do LED é limitada e, nessa condição, a corrente pode crescer rapidamente e danificar o componente.

\section{Aquisição de dados}

Registre cada medida com os dados necessários para justificar a incerteza escrita.

\subsection{Resistências medidas diretamente}

\small
\begin{longtable}{>{\raggedright\arraybackslash}p{1.7cm} >{\raggedright\arraybackslash}p{2cm} >{\raggedright\arraybackslash}p{1.8cm} >{\raggedright\arraybackslash}p{1.8cm} >{\raggedright\arraybackslash}p{3cm} >{\raggedright\arraybackslash}p{1.8cm} >{\raggedright\arraybackslash}p{2.3cm}}
\hline
Grandeza & Valor lido & Faixa usada & Resolução & Especificação do manual & Incerteza & Resultado final \\
\hline
\endfirsthead
\hline
Grandeza & Valor lido & Faixa usada & Resolução & Especificação do manual & Incerteza & Resultado final \\
\hline
\endhead
R0 &  &  &  &  &  &  \\
\hline
R1 &  &  &  &  &  &  \\
\hline
R2 &  &  &  &  &  &  \\
\hline
\end{longtable}
\normalsize

\subsection{Tensões medidas sobre os resistores}

\small
\begin{longtable}{>{\raggedright\arraybackslash}p{1.7cm} >{\raggedright\arraybackslash}p{2cm} >{\raggedright\arraybackslash}p{1.8cm} >{\raggedright\arraybackslash}p{1.8cm} >{\raggedright\arraybackslash}p{3cm} >{\raggedright\arraybackslash}p{1.8cm} >{\raggedright\arraybackslash}p{2.3cm}}
\hline
Grandeza & Valor lido & Faixa usada & Resolução & Especificação do manual & Incerteza & Resultado final \\
\hline
\endfirsthead
\hline
Grandeza & Valor lido & Faixa usada & Resolução & Especificação do manual & Incerteza & Resultado final \\
\hline
\endhead
V0 &  &  &  &  &  &  \\
\hline
V1 &  &  &  &  &  &  \\
\hline
V2 &  &  &  &  &  &  \\
\hline
\end{longtable}
\normalsize

\subsection{Corrente medida diretamente}

\small
\begin{tabular}{p{1.9cm} p{1.9cm} p{1.9cm} p{1.9cm} p{3.1cm} p{1.8cm} p{2.2cm}}
\hline
Grandeza & Valor lido & Faixa usada & Resolução & Especificação do manual & Incerteza & Resultado final \\
\hline
I\_medida &  &  &  &  &  &  \\
\hline
\end{tabular}
\normalsize

\section{Análise e interpretação dos dados}

Para três resistores em série:

\begin{equation}
R_{\mathrm{eq}} = R_0 + R_1 + R_2
\end{equation}

As duas formas de propagação ficam:

\begin{itemize}
  \item pela equação linear,
  \begin{equation}
  \Delta R_{\mathrm{eq}} = \Delta R_0 + \Delta R_1 + \Delta R_2
  \end{equation}
  \item pela equação em quadratura,
  \begin{equation}
  \Delta R_{\mathrm{eq}} = \sqrt{(\Delta R_0)^2 + (\Delta R_1)^2 + (\Delta R_2)^2}
  \end{equation}
\end{itemize}

Para estimar a corrente na parte resistiva do circuito:

\begin{equation}
I_{\mathrm{calc}} = \frac{V_0 + V_1 + V_2}{R_{\mathrm{eq}}}
\end{equation}

Usa-se a soma das quedas de tensão nos resistores, e não diretamente a tensão da fonte, porque o LED também participa da
distribuição de tensão.

Definindo a tensão total sobre a parte resistiva por

\begin{equation}
V_R = V_0 + V_1 + V_2
\end{equation}

as duas formas de propagação para a corrente podem ser escritas como:

\begin{itemize}
  \item forma linear,
  \begin{equation}
  \Delta I_{\mathrm{lin}} =
  \left| \frac{\partial I}{\partial V_R} \right| \Delta V_R +
  \left| \frac{\partial I}{\partial R_{\mathrm{eq}}} \right| \Delta R_{\mathrm{eq}}
  =
  \frac{\Delta V_R}{R_{\mathrm{eq}}} + \frac{V_R}{R_{\mathrm{eq}}^2}\Delta R_{\mathrm{eq}}
  \end{equation}
  \item forma em quadratura,
  \begin{equation}
  \Delta I_{\mathrm{quad}} =
  \sqrt{
  \left(\frac{\Delta V_R}{R_{\mathrm{eq}}}\right)^2 +
  \left(\frac{V_R}{R_{\mathrm{eq}}^2}\Delta R_{\mathrm{eq}}\right)^2
  }
  \end{equation}
\end{itemize}

Compare a corrente calculada com a corrente medida e verifique se a diferença é compatível com as incertezas experimentais.

\section{Perguntas conceituais e aplicadas}

\begin{enumerate}
  \item Por que uma medida de resistência não deve ser feita com o circuito energizado?
  \item Qual é a diferença prática entre ligar o multímetro em paralelo para medir tensão e em série para medir corrente?
  \item O que muda na interpretação do circuito quando há um LED em série com os resistores?
  \item Em que sentido a propagação linear é mais conservadora do que a propagação em quadratura?
  \item Se a corrente calculada e a corrente medida diferirem além das incertezas, quais causas experimentais podem explicar isso?
  \item Por que é importante registrar a faixa usada, a resolução e a especificação do manual junto com cada medida?
\end{enumerate}

\section{Simulador}

Para apoio ao relatório, você pode reconstruir o circuito no Tinkercad Circuits, disponível em
https://www.tinkercad.com/circuits, usando os mesmos elementos da prática: fonte, três resistores em série e LED.

O simulador não substitui a bancada na parte de incertezas instrumentais, porque não reproduz a especificação real do multímetro
usado no laboratório. Ainda assim, ele é útil para verificar qualitativamente as quedas de tensão nos resistores e comparar a
distribuição de tensão observada no circuito.

\section{Como organizar o relatório}

O relatório deve mostrar os resultados e como eles foram obtidos. Uma estrutura simples é suficiente:

\begin{enumerate}
  \item objetivo da prática;
  \item descrição curta do circuito e dos instrumentos usados;
  \item tabelas com as medidas de resistência, tensão e corrente;
  \item explicação de como a incerteza de cada medida direta foi obtida a partir do manual do multímetro;
  \item cálculo da resistência equivalente com propagação de incertezas pelos dois métodos;
  \item cálculo da corrente e comparação com a corrente medida diretamente;
  \item reconstrução do circuito no Tinkercad, para verificar qualitativamente as quedas de tensão nos resistores;
  \item discussão do papel do LED na interpretação do circuito;
  \item conclusão final, indicando o que foi aprendido e quais limitações apareceram.
\end{enumerate}

\paragraph{Dica prática.}
O aluno deve registrar o desenvolvimento das contas, ao menos uma vez de forma completa, e não apenas preencher a última linha com
o valor final.

\section{Discussão e conclusão}

Organize a discussão em torno das perguntas abaixo:

\begin{enumerate}
  \item As resistências medidas são consistentes entre si e com a escrita correta das incertezas?
  \item A resistência equivalente calculada é coerente com uma associação em série?
  \item A corrente calculada é compatível com a corrente medida diretamente?
  \item Que papel o LED desempenha na interpretação das quedas de tensão e dos limites da lei de Ohm neste circuito?
\end{enumerate}

Na conclusão, relacione técnica de medida, qualidade do registro experimental e interpretação física do circuito.

\section{Referências}

\begin{itemize}
  \item Material do portal da disciplina de Física Experimental IV.
\end{itemize}

\end{document}
